\subsection*{MT-M05. The stated Schauder quotient loses more than its zero evaluations see}
Lemma 9.3.2.5 on original PDF pp.~425--426 asserts a basis for the
quotient by the actual convolution ideal $\varphi*\mathcal S_E(G)$,
with coefficient functionals the Mellin jets at the zeros of
$\mathcal M\varphi$. We inspected the original p.~426 image: there
is no closure over this ideal. The assertion is false even when the
Mellin divisor consists of one simple zero. Thus the counterexample
does not rely on the zero-count defect in MT-M01.

Use $G=\mathbb R_{>0}$, trivial character and $u=\log t$; convolution
becomes ordinary additive convolution. Write
\[
 g_a(u)=(4\pi a)^{-1/2}e^{-u^2/(4a)},\qquad a>0.
\]
Choose $0<b<a$ and set $\varphi(e^u)=-g_a'(u)$,
$h(e^u)=-g_b'(u)$. Both functions belong to the exponential-Schwartz
space, and
\[
 \mathcal M\varphi(s)=s e^{as^2},\qquad
 \mathcal Mh(s)=s e^{bs^2}.
\]
The first function has just the simple zero $0$, lying in the interior
of the strip for $I=[-1,1]$. The only claimed coefficient functional is
$\delta_0$, and $\delta_0(h)=0$.
But $h\notin\varphi*\mathcal S_E$: if $h=\varphi*f$, multiplicativity
of the absolutely convergent Mellin transform gives
$\mathcal Mf(s)=e^{(b-a)s^2}$ for $s\ne0$, and then everywhere by
analyticity. On the imaginary axis this has modulus
$e^{(a-b)\tau^2}$. In contrast, $f(e^u)$ is integrable and therefore
$|\mathcal Mf(i\tau)|\le\int|f(e^u)|\,du$, a contradiction.
Thus the nonzero quotient class of $h$ has every asserted coefficient
zero, which is incompatible with uniqueness in a Schauder expansion.
The still simpler choice $\varphi=g_a$, $h=g_b$ has no Mellin zeros
and a nonzero quotient, so even the empty-divisor case fails.

\subsection*{MT-M06. The closed-ideal repair for every polynomial Gaussian divisor}
In logarithmic coordinates put
\[
 \mathcal E=\{f\in C^\infty(\mathbb R,\mathbb C):
   p_{N,j}(f):=\sup_x e^{N|x|}|f^{(j)}(x)|<\infty
   \text{ for every integer }N,j\ge0\}.
\]
This is $\mathcal S_E(\mathbb R_{>0})$ with its exponential-Schwartz
Fr\'echet topology. Polynomial factors in the usual Schwartz seminorms
are absorbed by increasing $N$. Let
\[
 P(s)=\prod_{\nu=1}^r(s-\rho_\nu)^{m_\nu},\quad
 d=\sum_\nu m_\nu,\qquad
 \varphi=P(-D)g_a,\quad D=d/dx.
\]
Distinct $\rho_\nu\in\mathbb C$, positive integral $m_\nu$, and
$a>0$ are arbitrary; $P=1,d=0$ is allowed. Then
\[
 \overline{\varphi*\mathcal E}
   =J_P:=\{h\in\mathcal E:
          (\mathcal Mh)^{(j)}(\rho_\nu)=0,
          \ 0\le j<m_\nu\},                         \tag{MT-M6}
\]
and the quotient by this \emph{closed} ideal is topologically
$\mathbb C^d$, with exactly the displayed Mellin jets as coordinates.
In particular its coordinate basis is absolute. This repairs the
finite-divisor case after making the closure operation explicit.

\begin{proof}
First, $g_a*\mathcal E$ is dense in $\mathcal E$. We give the full
argument, since density is distinct from the false surjectivity in
MT-M05. Convolution is continuous on $\mathcal E$: the inequality
$e^{N|x|}\le e^{N|x-y|}e^{N|y|}$ bounds each differentiated
convolution by an exponential supremum seminorm times an exponential
$L^1$ seminorm; the latter is bounded by a constant times
$p_{N+1,0}$.

Suppose a continuous linear functional $T$ vanishes on
$g_a*\mathcal E$. For every real $y$ and $\epsilon>0$ the function
$g_{a+\epsilon}(x-y)=g_a*g_\epsilon(\,\cdot-y)$ lies in this range.
It converges to $g_a(x-y)$ in every $p_{N,j}$ as
$\epsilon\downarrow0$, by uniform Gaussian estimates with the
parameter in a compact subinterval of $(0,\infty)$. Hence
$T(g_a(\,\cdot-y))=0$. The function
\[
 B(z)=T\bigl(e^{-x^2/(4a)}e^{zx/(2a)}\bigr)
\]
is entire: on every compact set of $z$, all its proposed complex
derivatives converge in every $p_{N,j}$ by Gaussian domination.
Completing the square gives $B(y)=0$ for real $y$, so the identity
theorem gives $B=0$.

For any Schwartz function $\psi$, Fourier inversion now gives
\[
 T(e^{-x^2/(4a)}\psi(x))
  =\frac1{2\pi}\int_{\mathbb R}\widehat\psi(\xi)
        B(2ai\xi)\,d\xi=0.
\]
This passage under $T$ is justified in the Fr\'echet space itself:
the $p_{N,j}$ seminorm of
$e^{-x^2/(4a)}e^{i\xi x}$ is at most
$C_{N,j,a}(1+|\xi|)^j$, whereas $\widehat\psi$ is rapidly
decreasing. For $h\in C_c^\infty$, take
$\psi=e^{x^2/(4a)}h$, which is again compactly supported and smooth.
Then $T(h)=0$. Finally $C_c^\infty$ is dense in $\mathcal E$:
multiply by a smooth cutoff equal to one on $[-R,R]$ and supported
in $[-2R,2R]$; the product rule and the bound from $p_{N+1,j}$ make
every $p_{N,j}$ of the discarded tail tend to zero. Thus $T=0$.
Hahn--Banach separation proves the asserted density.

Second, the range of $P(-D)$ on $\mathcal E$ is exactly $J_P$.
Integration by parts gives $\mathcal M(-Df)(s)=s\mathcal Mf(s)$,
so inclusion in $J_P$ follows. For the converse, if
$\mathcal Mh(\rho)=0$, define
\[
 f(x)=e^{-\rho x}\int_{-\infty}^x e^{\rho t}h(t)\,dt
     =-e^{-\rho x}\int_x^\infty e^{\rho t}h(t)\,dt.
\]
The equality is precisely the vanishing of the Mellin value.
The first expression controls the negative tail and the second the
positive tail. For any prescribed $N$, using an exponential bound
on $h$ with exponent larger than $N+|\operatorname{Re}\rho|+1$
proves exponential decay of $f$ on both tails; the equation
$f'=h-\rho f$ then gives the same bounds for every derivative.
Therefore $f\in\mathcal E$, and $(D+\rho)f=h$.
Equivalently $-f$ solves $(-D-\rho)(-f)=h$, and its Mellin transform
is $\mathcal Mh(s)/(s-\rho)$. Repeating this division through all
roots, including multiplicities, proves $P(-D)\mathcal E=J_P$.

The Mellin jet functionals are continuous, so $J_P$ is closed.
Since $g_a*\mathcal E$ is dense and $P(-D)$ is continuous,
$P(-D)(g_a*\mathcal E)$ is dense in $P(-D)\mathcal E=J_P$.
Differentiation commutes with convolution, identifying that range
with $\varphi*\mathcal E$. This proves (MT-M6).

For $d=0$ the quotient is zero and no interpolation is required.
For $d>0$, it remains to verify that no coordinate is missing. Given arbitrary
$d$ target jets, ordinary polynomial Hermite interpolation constructs
a polynomial $Q$ of degree below $d$ for which
$Q(s)e^{bs^2}$ has exactly those jets, with any fixed $b>0$.
Indeed multiplication by the nowhere-zero entire function $e^{bs^2}$
is invertible on each finite jet space, so interpolate the resulting
jets of $Q$. The function $Q(-D)g_b$ belongs to $\mathcal E$ and
realizes them. Choosing the targets to be the coordinate vectors
gives a continuous linear section from the finite-dimensional space
$\mathbb C^d$ to $\mathcal E$. The jet map consequently induces
a continuous isomorphism of the quotient with $\mathbb C^d$,
with continuous inverse. Finite-dimensional coordinate expansions are
absolute, completing the proof.
\end{proof}

For $P=s$, MT-M05 therefore exhibits a proper, nonclosed convolution
ideal whose closure is $\ker\delta_0$. The closed quotient really
does have the one coordinate the draft intended. For $P=1$ the
principal ideal is dense but proper, and the closed quotient is zero.
This is a concrete repair rather than a certification of the general
infinite-divisor statement. For arithmetic Hecke divisors the remaining
gate is controlled infinite interpolation and the claimed absolute
basis/nuclear estimates on the correct separated quotient. A local
$O(\log T)$ zero count alone does not prove those estimates.
In cote~1 trace language, Hausdorff completion and the identification
of the spectral quotient precede taking its nuclear trace.
